\chapter{Governança: Estoque e Matriz de Acesso}

A sustentabilidade financeira e a segurança jurídica da rede IntraClinica residem no controle rigoroso sobre os insumos e sobre quem manipula a informação.

\section{O Guardião de Estoque Inteligente}

O estoque no IntraClinica não é um módulo isolado, mas uma consequência do atendimento.
\begin{itemize}
    \item \textbf{Receita de Procedimento:} Para cada ato médico (ex: aplicação de injetável, pequena cirurgia), existe uma ``receita'' cadastrada no sistema.
    \item \textbf{Baixa Atômica:} Ao finalizar o procedimento no prontuário, o sistema automaticamente propõe a baixa dos insumos associados. Se uma sutura foi realizada, o kit de sutura e os fios são deduzidos do inventário instantaneamente.
    \item \textbf{Alertas de Anomalia:} A inteligência local monitora desvios estatísticos. Se o consumo de um determinado profissional ou turno diverge do padrão para aquele procedimento, a Diretoria Médica recebe um alerta silencioso de auditoria.
\end{itemize}

\section{Gestão de Identidade e Acesso (IAM)}

Para viabilizar a curadoria de médicos parceiros, implementamos uma matriz de acesso baseada no princípio do menor privilégio.

\begin{table}[h!]
\centering
\caption{Matriz de Funções e Escopos}
\vspace{0.3cm}
\begin{tabular}{@{}ll@{}}
\toprule
\textbf{Papel (Role)} & \textbf{Responsabilidade e Acesso} \\ \midrule
Diretoria Médica & Gestão total, auditoria e definição de protocolos. \\
Médico Especialista & Acesso pleno ao prontuário e atendimento clínico. \\
Secretariado & Foco em agendamento, check-in e dados cadastrais. \\
Gestor de Estoque & Controle de suprimentos, validades e custos. \\ \bottomrule
\end{tabular}
\end{table}

Essa estrutura garante que a intimidade do paciente (contida no prontuário médico) seja acessada apenas por profissionais autorizados, cumprindo rigorosamente as normativas do CFM e da LGPD.
